% =============================================================================
% Online Appendix to Section 6 (Economic Opportunity): construction details and further results.
% Numbers: paper/results/eo.json (scripts/paper/fig_eo.R) and eo_transforms.json
% (scripts/paper/fig_eo_transforms.R), final rerun of 2026-10-12. Statistics over rankable cells
% (config/rankability.yml) unless stated.
% =============================================================================
\section{Economic Opportunity: construction details and further results}\label{oa:eo}

\subsection{Construction details}

\paragraph{Trade.} Revealed comparative advantage complements market share: a small economy can be
highly specialised in a basket without a large share of the world market, and specialisation is
what the complexity literature reads as revealed capability. Basket export value is a second
measure of scale from a source with different coverage and year. Feasibility is the forward-looking
element: a country with no exports of a basket can score well if it exports products that require
similar capabilities. An economy that does not report to Comtrade is scored on its two Atlas
components and the world-market term. The net-import term of the ES category is not used in EO.

\paragraph{Production.} Absolute growth rewards large producers whose output expands by the most
units, relative growth rewards fast growth from a small base, and the composite balances the two;
level and change each enter the category twice, directly and through the composite. In the mineral
block, tonnes of different minerals cannot be added, so each mineral enters as a share of world
supply weighted by the technology's share of the mineral's demand; the horizon level is divided by
base-year world supply, so the change is a change in quantity, not in market share. A country the
IEA supply table does not name for a mineral counts as zero there. Uranium mine output from the Red
Book is scored like oil, gas and coal.

\paragraph{Technology Demand.} Growth in the World Energy Outlook composite is absolute because a
relative measure ranked across fuels would put hydrogen's near-zero base at the top; the value is a
property of the technology, the same at every stage and for every economy. The BCG addressable
market is ranked across the 26 cells it sizes and copied to every economy by the global-average
rule. Coverage is uneven: the electric grid has no input in any source and so no Technology Demand
category, electric vehicles have it only downstream, downstream batteries lack it because BCG does
not size that cell, and geothermal is scored on the BCG entry alone.

\paragraph{Technological Readiness.} The bell rises over four steps and falls over five, so an
item at TRL 9 keeps a level score of $\cos(3\pi/10)\approx0.59$ and one at TRL 3 scores
$\cos(3\pi/8)\approx0.38$; the concave momentum term gives the first step of progress the largest
increment. Of the 253 mapped items, 9 map to more than one cell; 98 carry a post-TRL-9 adoption
stage, 43 per cent sit at TRL 9 or above, and the 26 items that are commercial in many markets score
zero on level. Of the items with a 2020 reading, 54 per cent did not move. Ten of the 33 headline
cells take their technology's value because no item maps to the stage; the index differs across the
three stages of ten of the eleven technologies (wind keeps one value), with a median range across a
technology's stages of 0.16. The parameters (bounds 2 and 11, peak 6, momentum weight 0.3, cap of
three steps, exponent 1/2) are the theme builder's defaults.

\paragraph{Overcapacity.} Of the four midstream technologies the BloombergNEF supply-chain data
cover, solar ($\rho=2.3$, index 0.20), batteries and wind ($\rho=3.0$, 0.13 each) sit on the steep
part of the curve, and green hydrogen ($\rho=18.2$) is effectively at zero.

\paragraph{Cost Competitiveness.} Wages are converted at market exchange rates rather than
purchasing-power parities because the price-level difference that PPP conversion removes is the
cost advantage being measured. The lending rate is winsorised at its 5th and 95th cross-country
percentiles before min-max rescaling; producer-price inflation enters as a rate because each
country's price index has its own base year. Stale observations and gaps are filled from regional
and income-group peers. The Input Cost Index carries the category wherever the LCOE and equipment
cost components are absent and half of it where one is present; counting each source once avoids
giving BNEF three of four votes wherever it covers a market.

\subsection{Further results}

\paragraph{Rankability.} Of the 2{,}830 cells the rule sets aside, 2{,}112 carry more than half of
their weight on imputed components, 74 belong to territories and 644 fail on both counts. The rule
bites hardest where country-specific sources are thinnest: 70.2 per cent of upstream cells are
rankable, against 55.7 per cent midstream and 54.9 per cent downstream. In Batteries midstream 42
of the 54 rankable economies have an observed Foreign Dependency or Investment input, in Green
Hydrogen midstream 30 of 40. The mean imputed weight share of rankable cells is 0.29. Ranked among
all scored economies of their pair, cells excluded for imputation alone have a mean percentile of
0.43 (63.4 per cent below the median) and territory cells 0.42, against 0.55 for rankable cells.
Both decompositions reconcile exactly (largest gaps 0 and $1.1\times10^{-16}$); 45.3 per cent of
category scores contain an imputed component.

\paragraph{Distribution.} Over all scored cells the midstream median is 0.454 and its lower
quartile 0.328, because the zero rule places economies without recorded capacity or investment at
the bottom of the clean-technology midstream pairs, and the rule removes most of them. Four of the
five clean-technology midstream pairs (wind, electric vehicles, solar, batteries) have the lowest
midstream medians (0.351 to 0.426). The widest boxes owe their width to one category each: in EV
midstream (interquartile range 0.144) Foreign Dependency, observed for 97 of 145 rankable
economies, accounts for 46 per cent of the variance (0.053 without it); in Green Hydrogen
downstream (0.122) Consumption accounts for 94 per cent; in Wind downstream (0.117) Investment
accounts for 46 per cent.

\paragraph{Variance and persistence.} Technology effects alone explain 27.8 per cent of the
variance and stage effects almost nothing (0.01 per cent), so a pair's level is set by its
composition rather than its stage. Over all scored cells the pair share is 71.5 per cent. Across
the 159 economies rankable at all three stages, the Spearman correlation of mean within-pair
percentiles is 0.86 between upstream and midstream, 0.70 between upstream and downstream and 0.79
between midstream and downstream. Part of this common factor is the Input Cost Index, a
country-level index whose stage variation comes only from the wage activity and the factor shares:
across the 178 economies with a rankable cell, mean EO percentile and mean Cost Competitiveness
percentile correlate at 0.77 (0.82 across the 57 coalition members).

\paragraph{Drivers.} Level shares: upstream Cost Competitiveness 21.9 per cent, Technology Demand
17.3, Production 16.5 and Energy Prices 16.2; midstream Cost Competitiveness 29.2, Technology Demand
24.7 and Trade 24.6; downstream Cost Competitiveness 27.5, Technology Demand 19.7 and Technological
Readiness 18.0. Within-pair variance shares: upstream Trade 26.9 and Production 22.4; midstream
Trade 25.5 and Foreign Dependency 22.2; downstream Consumption 23.8 and Investment 22.2. Investment
contributes 0.3, 1.2 and 4.2 per cent of the level by stage but 8.3, 12.7 and 22.2 per cent of the
within-pair variance, because under the zero rule the few economies with recorded investment stand
far above the majority without it (in 9 of the 12 pairs that carry the category). The upstream
constants are led by price volatility (15.9 per cent of effective weight), the TRL Index (12.5) and
the addressable market (10.1). Before the rankability rule the observed country-specific share is
32.3 per cent and the imputed share 30.1 per cent.

\paragraph{Leaders.} Chad combines low recorded wages with a 2017 lending rate of 2.45 per cent, the
joint fifth lowest of the 144 rates in the cost theme; the runner-up in its three pairs is Cameroon,
with the same rate (2018), and Chad's margins are 0.0020, 0.0005 and 0.0011. Canada leads nuclear
upstream (0.65) and China electric-grid upstream (0.71, ahead of Chile at 0.65), the pair built on
the copper and aluminium basket. The leader's margin is below 0.01 in seven pairs; the United States
has the largest non-Chinese margin (0.153 over Indonesia in geothermal downstream). Coalition
members hold 69.1 per cent of the 330 top-ten places; China is in the top ten of 30 pairs and the
United States of 20. The United States ranks in the bottom half of one pair only, 109th of 141 in
electric-grid downstream, a cost-only pair; in the mineral-based upstream pairs it ranks between
third (wind) and 62nd (green hydrogen). China's one low rank, 98th of 104 in oil downstream, comes
from Consumption and the growth component of Production.

\paragraph{The coalition.} Among core members the United States (0.89), Portugal (0.84), Poland
(0.84), the Czech Republic (0.81) and Spain (0.80) rank highest and the United Kingdom (0.57),
Romania (0.52), Australia (0.46), New Zealand (0.34) and Ukraine (0.11) lowest; among partners,
Thailand (0.82), Indonesia (0.80), Malaysia (0.79), India (0.78) and Saudi Arabia (0.66) rank highest
and the Dominican Republic (0.26), Brazil (0.23), Argentina (0.19), the Democratic Republic of the
Congo (0.13) and Ghana (0.10) lowest. T\"urkiye, Brazil, Argentina and Ghana record nominal lending
rates of 39.0, 46.0, 59.7 and 27.0 per cent, at the winsorisation ceiling; with the Cost
Competitiveness contribution set to its pair mean their mean percentiles would be 0.73, 0.61, 0.45
and 0.39 rather than 0.30, 0.23, 0.19 and 0.10, and Ukraine's (rate 18.9 per cent) 0.50 rather than
0.11; six members gain more than 0.2. At the top the same change would lower Portugal from 0.84 to
0.62, Thailand from 0.82 to 0.64 and Ireland from 0.71 to 0.51, while the United States barely
moves. The BNEF regional composite of Technology Demand is observed for at most twelve economies
and the global-average rule gives the others their mean, so a modelled economy below that mean
ranks below many unmodelled ones: Australia's value is below the mean in all 18 pairs in which it is
observed, the United Kingdom's in 14 of 18 and France's in 15 of 21; with Technology Demand at its
pair mean their mean percentiles would be 0.52, 0.65 and 0.80 rather than 0.46, 0.57 and 0.76.

Of the core members' 1{,}023 scored cells 1{,}002 are rankable, against 716 of the partners' 858.
Gabon has only 9 rankable cells, Oman 21, Saudi Arabia and Singapore 22 each, and Morocco, Namibia
and Costa Rica 23. The mean imputed weight share of rankable cells is 0.20 for the core, 0.27 for
the partners, 0.33 for other economies and zero for the United States, partly because the BCG
addressable market is recorded for the United States alone. Within pairs the median Spearman
correlation between $\EO_{its}$ and the imputed share is $-0.28$.

\paragraph{Illustrative pairs.} Oil upstream is led by the United States (0.659), China (0.644),
Canada (0.637), Saudi Arabia (0.622) and Norway (0.607); Saudi Arabia, Iraq, Iran, Kuwait and the
United Arab Emirates share the smallest Cost Competitiveness contribution (0.095, an imputed value)
and owe their places to production and trade. In Batteries midstream China (0.771) leads the United
States (0.633) on Foreign Dependency (0.167 against 0.106) and Cost Competitiveness (0.142 against
0.073), while the United States' Investment contribution (0.143) exceeds China's (0.134); nine of
the ten are coalition members. In Solar midstream China (0.788) is followed by India (0.640) and
the United States (0.635), then Viet Nam, Thailand, Indonesia, Malaysia and the Philippines, with
France and Germany eighth and ninth. In EV downstream emerging economies take eight of the ten
places: China (0.755), Viet Nam (0.731), Thailand (0.705) and T\"urkiye (0.702) lead, because the
EV-market composite, half of whose indices are growth rates (equation~\eqref{eq:eo_ev}), rewards
markets growing fast from a small base; T\"urkiye, Viet Nam and the Philippines receive larger
Technology Demand contributions than China.
